% ===========================================================================
% ALGORITHM / ALGPSEUDOCODE CHEAT SHEET
% ---------------------------------------------------------------------------
% Basic statement:
%   \State do something
%
% Assignment:
%   \State $x \gets y + z$
%
% Comments (inline):
%   \State $x \gets 0$ \Comment{this is a comment}
%
% If / ElsIf / Else:
%   \If{$x > 0$}
%     \State ...
%   \ElsIf{$x = 0$}
%     \State ...
%   \Else
%     \State ...
%   \EndIf
%
% For loop:
%   \For{$i = 1$ \textbf{to} $N$}
%     \State ...
%   \EndFor
%
% For-each loop:
%   \ForAll{$x \in S$}
%     \State ...
%   \EndFor
%
% While loop:
%   \While{$x > 0$}
%     \State ...
%   \EndWhile
%
% Repeat-Until:
%   \Repeat
%     \State ...
%   \Until{condition}
%
% Function definition:
%   \Function{MyFunc}{$a, b$}
%     \State \Return $a + b$
%   \EndFunction
%   \State $z \gets \Call{MyFunc}{x, y}$
%
% Require / Ensure (preconditions / postconditions):
%   \Require input is valid
%   \Ensure output satisfies property
%
% Line numbers: enabled by \begin{algorithmic}[1]  (step size = 1)
%
% Cross-reference an algorithm:
%   \label{alg:foo}  inside \begin{algorithm}
%   refer with \cref{alg:foo}  (needs cleveref) or Algorithm~\ref{alg:foo}
% ===========================================================================

\documentclass[11pt]{article}
\usepackage[margin=1in]{geometry}
\usepackage{algorithm}
\usepackage{algpseudocode}
\usepackage{cancel}
\usepackage{graphicx}
\usepackage{float}
\usepackage{amsmath, amsfonts, amssymb, esint, multicol}
\usepackage{hyperref}
\usepackage{cleveref}
\usepackage[framemethod=tikz]{mdframed}
\usepackage{xcolor}
\usepackage{tikz}
\usetikzlibrary{arrows.meta,positioning}
\newcommand*\circled[1]{\tikz[baseline=(char.base)]{
            \node[shape=circle,draw,inner sep=2pt] (char) {#1};}}
% Define custom color and box environment
\definecolor{mycolor}{rgb}{0.122, 0.435, 0.698}
\newmdenv[
    innerlinewidth=0.5pt,
    roundcorner=4pt,
    linecolor=mycolor,
    innerleftmargin=10pt,
    innerrightmargin=10pt,
    innertopmargin=4pt,
    innerbottommargin=3pt
]{mybox}

% MACROS
\newcommand{\R}{{\mathbb R}}
\newcommand{\V}{\vec}
\newcommand{\resetcount}{{\setcounter{equation}{0}
}}
\DeclareMathOperator*{\argmin}{arg\,min}

\usepackage{parskip}
\newcommand{\volume}{\makebox[1pt][l]{$-$}V}

\usepackage{listings}
\lstset{ 
  breaklines=true,
  breakatwhitespace=true
}

\usepackage{fancyhdr}
\usepackage{titlesec}

% Title page formatting
\title{\vspace{2in}\textbf{NUEN 625 Project}\\\vspace{0.5in}\large Neutron Transport Theory}
\author{Dylan Frank\\Texas A\&M University}
\date{Due: Friday, March 24, 2026}

\begin{document}

% Cover page
\maketitle
\thispagestyle{empty}
\newpage

% Start page numbering
\setcounter{page}{1}
\section*{Problem 9}

Show that if one assumes that
\begin{equation*}
  \psi(-1) = (\phi-3J)/2
\end{equation*}  

and further that

\begin{equation*}
  \psi(\mu_m) = (\phi-3J\mu_m)/2,\quad m=1,\dots, N,
\end{equation*}

then the weighted-diamond difference approximation implies that

\begin{equation*}
  \psi(\mu_{m+1/2}) = (\phi-3J\mu_{m+1/2})/2,\quad m=0,\dots, N,
\end{equation*}

\begin{mybox}
  We proceed by induction on $m$. The base case $m=0$ follows immediately: since $\mu_{1/2} = -1$, the assumption $\psi(-1) = (\phi-3J)/2$ gives $\psi_{1/2} = (\phi - 3J\mu_{1/2})/2$. For the inductive step, assume $\psi_{m-1/2} = (\phi-3J\mu_{m-1/2})/2$ and note that the weighted-diamond difference closure is
  \begin{equation*}
    \psi_m = \beta_m \psi_{m+1/2} + (1-\beta_m)\psi_{m-1/2}\quad .
  \end{equation*}
  
  Rearranging to isolate for $\psi_{m+1/2}$ and substituting the ansatze $\psi_m = (\phi-3J\mu_m)/2$, $\psi_{m-1/2} = (\phi-3J\mu_{m-1/2})/2$ gives us
  \begin{align}
    \psi_{m+1/2} &= \frac{1}{\beta_m} [\psi_m + (\beta_m -1)\psi_{m-1/2}]\nonumber \\
    &=\frac{1}{\beta_m} \left[\frac{\phi-3J\mu_m}{2} + (\beta_m-1)\frac{\phi-3J\mu_{m-1/2}}{2}\right], \nonumber\\
    &= \frac{1}{\beta_m} \frac{3J\mu_{m-1/2}-3J\mu_m}{2} + \frac{\phi-3J\mu_{m-1/2}}{2}\quad. \label{eq:9-1}
  \end{align}

  Using the definition $\beta_m = (\mu_m-\mu_{m-1/2}) / (\mu_{m+1/2}-\mu_{m-1/2})$, Eq.\eqref{eq:9-1} becomes
  \begin{align*}
    \psi_{m+1/2} &= \frac{\mu_{m+1/2}-\mu_{m-1/2}}{\mu_m-\mu_{m-1/2}} \cdot \frac{3J\mu_{m-1/2}-3J\mu_m}{2} + \frac{\phi-3J\mu_{m-1/2}}{2} \\
    &= \frac{-3J(\mu_m-\mu_{m-1/2})(\mu_{m+1/2}-\mu_{m-1/2})}{2(\mu_m-\mu_{m-1/2})} + \frac{\phi-3J\mu_{m-1/2}}{2} \\
    &= \frac{\phi - 3J\mu_{m-1/2} - 3J(\mu_{m+1/2}-\mu_{m-1/2})}{2} \\
    &= \frac{\phi - 3J\mu_{m+1/2}}{2}\quad.
  \end{align*}


\end{mybox}



\section*{Problem 10}
Starting with the spatially-analytic S$_n$ equations, assume that
\begin{equation*}
  \psi(\mu_j) = (\phi + 3J\mu_j )/2 , \quad j = 1/2, 1, 3/2, ...., N, N + 1/2,
\end{equation*}
and show that taking the zero’th and first angular moments of the S$_n$ equations using the quadrature formula yields the analytic P$_1$ equations. Remember that any standard Sn quadrature set must integrate 1, $\mu^2$, and all odd functions of $\mu$ (the latter property is due to quadrature point symmetry about $\mu$ = 0). Also remember that the $\alpha$-coefficients are positive and even in $\mu$. You will need to take the symmetry into account to show that certain sums are equal to zero. This requires some thought, but it is a good exercise. The best way to determine if the spatially analytic Sn equations preserve the thick diffusion limit is to do an asymptotic analysis.

\begin{mybox}
  The S$_N$ equation analytic in space (and discrete in angle) for 1-D spherical coordinates is
  \begin{equation}
    \frac{\mu_m}{r^2}\frac{\partial[r^2\psi_m]}{\partial r} + \frac{[(\alpha_{m+1/2}\psi_{m+1/2} -\alpha_{m-1/2}\psi_{m-1/2})]}{r w_m} + \sigma_{t}\psi_m = Q_m \label{eq:10-1},
  \end{equation}
  where
  \begin{equation}
    Q_m = \sum_{k=0}^K \frac{2k+1}{2}(\sigma_k\phi_k + q_k)P_k(\mu_m)
  \end{equation}


  Evaluating the 0th moment (multiplying Eq.\eqref{eq:10-1} by $w_m$ and summing over $m\le N$), we first note that the second term (streaming) cancels out to 0 since $\alpha_{1/2} = \alpha_{N+1/2}$. The spatial term evaluates as,
  \begin{equation}
    \sum_{m=1}^N w_m\frac{\mu_m}{r^2}\frac{\partial[r^2\psi_m]}{\partial r} = \frac{1}{r^2}\frac{\partial[r^2 J]}{\partial r}\quad .
  \end{equation}
  
  The collision term evaluates to $\sigma_t\phi$, and the 0th moment of the scattering source is $\sigma_{s,0}\phi$. Combining terms and moving scattering to the left-hand side gives the 0th moment conservation equation,
  \begin{equation}
    \frac{\partial[r^2 J]}{\partial r} + \underbrace{(\sigma_t - \sigma_{s,0})}_{\sigma_a}\phi = q_0\,, \quad q_0 = \sum_m w_m q(\mu_m)\,.\label{eq:10-2}
  \end{equation}

  Before evaluating the 1st moment we establish two quadrature identities. Since the quadrature weights sum to 2 (one biradian), the set exactly integrates $\mu^2$ over $[-1,1]$:
  \begin{equation*}
    \sum_{m=1}^N w_m \mu_m^2 = \int_{-1}^{1} \mu^2\,d\mu = \left[\frac{\mu^3}{3}\right]_{-1}^{1} = \frac{2}{3}\,.
  \end{equation*}
  Because the abscissae are symmetric about $\mu=0$ (i.e.\ $\mu_{N+1-m} = -\mu_m$ with equal weights), every odd-power sum vanishes:
  \begin{equation*}
    \sum_{m=1}^N w_m \mu_m^3 = \int_{-1}^{1} \mu^3\,d\mu = 0\,.
  \end{equation*}

  Next, we evaluate the 1st moment equation (multiplying Eq.\eqref{eq:10-1} by $\mu_m w_m$ and summing over $m\le N$). For the spatial term we evaluate using the given ansatz and the identities above:
  \begin{align*}
    \sum_{m=1}^N \frac{w_m}{2r^2}\frac{\partial[r^2(\phi\mu_m^2 + 3J\mu_m^3)]}{\partial r} = \frac{1}{3r^2}\cdot \frac{\partial[r^2\phi]}{\partial r}\quad .
  \end{align*}

  For the angular term, we substitute the ansatz $\psi_{m+1/2} = (\phi + 3J\mu_{m+1/2})/2$ and split into two contributions:
  \begin{align*}
    T \equiv \frac{1}{r}\sum_{m=1}^N \mu_m\left(\alpha_{m+1/2}\psi_{m+1/2} - \alpha_{m-1/2}\psi_{m-1/2}\right)
    &= \underbrace{\frac{\phi}{2r}\sum_{m=1}^N \mu_m\!\left(\alpha_{m+1/2}-\alpha_{m-1/2}\right)}_{T_\phi}\\
    &+ \underbrace{\frac{3J}{2r}\sum_{m=1}^N \mu_m\!\left(\alpha_{m+1/2}\mu_{m+1/2}-\alpha_{m-1/2}\mu_{m-1/2}\right)}_{T_J}.
  \end{align*}

  For $T_\phi$: the $\alpha$-recursion gives $\alpha_{m+1/2} - \alpha_{m-1/2} = -2\mu_m w_m$, so
  \begin{equation*}
    T_\phi = \frac{\phi}{2r}\sum_{m=1}^N \mu_m(-2\mu_m w_m) = -\frac{\phi}{r}\sum_{m=1}^N w_m\mu_m^2 = -\frac{\phi}{r}\cdot\frac{2}{3} = -\frac{2\phi}{3r}\,.
  \end{equation*}

  For $T_J$: since $\alpha$ is even in $\mu$ ($\alpha_{N-m+1/2}=\alpha_{m+1/2}$) and $\mu$ is \emph{odd} ($\mu_{N+1-m}=-\mu_m$, $\mu_{N-m+1/2}=-\mu_{m+1/2}$), the contributions from symmetric pairs $(m,\,N+1-m)$ are equal and opposite, giving $T_J = 0$.

  Therefore $T = -2\phi/(3r)$, and the 1st-moment equation is
  \begin{equation*}
    \frac{1}{3r^2}\frac{\partial[r^2\phi]}{\partial r} - \frac{2\phi}{3r} + \sigma_t J = q_1\,.
  \end{equation*}
  Expanding the spatial term, $\tfrac{1}{3r^2}\partial_r(r^2\phi) = \tfrac{1}{3}\partial_r\phi + \tfrac{2\phi}{3r}$, the geometric terms cancel and we obtain the analytic P$_1$ current equation,
  \begin{equation}
    \frac{1}{3}\frac{\partial\phi}{\partial r} + \sigma_t J = q_1\,. \label{eq:10-3}
  \end{equation}
  Defining the diffusion coefficient $D \equiv \tfrac{1}{3\sigma_t}$, Eq.\eqref{eq:10-3} with $q_1=0$ becomes Fick's law,
  \begin{equation}
    J = -D\frac{\partial\phi}{\partial r}\,. \label{eq:fick}
  \end{equation}

\end{mybox}

\section*{Problem 11}
Derive a nearly-consistent diffusion equation for use in DSA.
\begin{mybox}
  We begin by discretizing the P$_1$ balance equation in Eq.\eqref{eq:10-2} for cell $i$ assuming a constant cross-section and source density within each cell,
  \begin{equation}
    A_{i+1/2}J_{i+1/2} - A_{i-1/2}J_{i-1/2} + V_i\sigma_{a,i} \phi_i = V_i q_{0,i} \quad,\label{eq:11-1}
  \end{equation}

  where $A_{i\pm1/2} = 4\pi r_{i\pm1/2}^2$ and $V_i = \frac{4}{3}\pi\left(r_{i+1/2}^3 - r_{i-1/2}^3 \right)$.
  
  The balance equation for cell $i+1$ is,
  \begin{equation}
    A_{i+3/2}J_{i+3/2} - A_{i+1/2}J_{i+1/2} + V_{i+1}\sigma_{a,i+1} \phi_{i+1} = V_{i+1} q_{0,i+1} \quad.\label{eq:11-2}
  \end{equation}

  Adding Eqs.\eqref{eq:11-1} and \eqref{eq:11-2} and using the approximation
  \begin{equation}
    A_{i+3/2}J_{i+3/2} - A_{i-1/2} J_{i-1/2} \approx 2(A_{i+1}J_{i+1} - A_iJ_i)\,,\label{eq:11-approx}
  \end{equation}
  where $A_i = (A_{i+1/2}+A_{i-1/2})/2$ is the arithmetic-mean face area and $J_i$ is the volume-averaged current, eliminates the half-index cell-edge currents:
  \begin{equation}
    2(A_{i+1}J_{i+1} - A_i J_i) + V_{i}\sigma_{a,i} \phi_{i} + V_{i+1}\sigma_{a,i+1} \phi_{i+1} = V_{i} q_{0,i}+ V_{i+1} q_{0,i+1}\quad .\label{eq:11-3}
  \end{equation}

  We now discretize the 1st-moment (Fick's law) equation using diamond differencing. This yields
  \begin{equation}
    J_i = -D_i \frac{A_i}{V_i}\left(\phi_{i+1/2} - \phi_{i-1/2}\right), \quad D_i = \frac{1}{3\sigma_{t,i}}\quad,\label{eq:11-4}
  \end{equation}
  where the $A_i/V_i$ factor carries the spherical-geometry correction. Substituting Eq.\eqref{eq:11-4} into Eq.\eqref{eq:11-3}, defining $\widetilde{D}_i \equiv D_i A_i^2/V_i$, and using the diamond approximation $\phi_i = (\phi_{i-1/2}+\phi_{i+1/2})/2$ yields the interior tridiagonal equation for $i = 1,\ldots,N-1$:
  \begin{align*}
    \underbrace{\left(-2\widetilde{D}_i + \tfrac{V_i\sigma_{a,i}}{2}\right)}_{\displaystyle a_i}\phi_{i-1/2}
    + \underbrace{\left(2\widetilde{D}_i + 2\widetilde{D}_{i+1} + \tfrac{V_i\sigma_{a,i}}{2} + \tfrac{V_{i+1}\sigma_{a,i+1}}{2}\right)}_{\displaystyle b_i}\phi_{i+1/2}
    \\+ \underbrace{\left(-2\widetilde{D}_{i+1} + \tfrac{V_{i+1}\sigma_{a,i+1}}{2}\right)}_{\displaystyle c_i}\phi_{i+3/2}
    = V_i q_{0,i} + V_{i+1} q_{0,i+1}\quad .\label{eq:11-5}
  \end{align*}

  \textbf{Boundary conditions.} At the inner boundary ($r=0$), symmetry requires $J_{1/2}=0$. Using only the $i=1$ cell balance Eq.\eqref{eq:11-1} with $J_{1/2}=0$ and substituting Eq.\eqref{eq:11-4} and $\phi_1 = (\phi_{1/2}+\phi_{3/2})/2$:
  \begin{equation}
    \underbrace{\left(2\widetilde{D}_1 + \tfrac{V_1\sigma_{a,1}}{2}\right)}_{\text{diagonal}}\phi_{1/2}
    + \underbrace{\left(-2\widetilde{D}_1 + \tfrac{V_1\sigma_{a,1}}{2}\right)}_{\text{off-diagonal}}\phi_{3/2}
    = V_1 q_{0,1}\,.\label{eq:11-bc-inner}
  \end{equation}

  At the outer boundary ($r=R$), a vacuum condition is imposed via the Marshak condition: the incoming partial current vanishes, giving $J_{N+1/2} = \langle\mu\rangle\,\phi_{N+1/2}$ where $\langle\mu\rangle = \sum_{\mu_m>0}\mu_m w_m$. Applying this to the $i=N$ cell balance Eq.\eqref{eq:11-1} and substituting Eq.\eqref{eq:11-4}:
  \begin{equation}
    \underbrace{\left(-2\widetilde{D}_N + \tfrac{V_N\sigma_{a,N}}{2}\right)}_{\text{off-diagonal}}\phi_{N-1/2}
    + \underbrace{\left(2A_{N+1/2}\langle\mu\rangle + 2\widetilde{D}_N + \tfrac{V_N\sigma_{a,N}}{2}\right)}_{\text{diagonal}}\phi_{N+1/2}
    = V_N q_{0,N}\,.\label{eq:11-bc-outer}
  \end{equation}

\end{mybox}

% OUTLINE FOR PROBLEM 11: Derive a nearly-consistent DSA diffusion equation
% -------------------------------------------------------------------------
%
% Goal: Starting from the two discrete P1 equations (balance + Fick's law),
%       eliminate the cell-edge currents to obtain a tridiagonal system for
%       the cell-edge scalar flux.
%
% Step 1 — Discretize the P1 balance equation (Eq. from Problem 10) for
%           cell i using diamond differencing:
%             A_{i+1/2} J_{i+1/2} - A_{i-1/2} J_{i-1/2} + sigma_a,i V_i phi_i = q_0,i V_i
%           where A_{i+/-1/2} = 4*pi*r_{i+/-1/2}^2 are the cell-face areas
%           and V_i = (4/3)*pi*(r_{i+1/2}^3 - r_{i-1/2}^3) is the cell volume.
%           The cell-center flux phi_i is approximated as the average of the two
%           cell-edge values: phi_i = (phi_{i-1/2} + phi_{i+1/2})/2.
%
% Step 2 — Write the same balance equation for cell i+1, then ADD the two
%           equations together. Use the approximation
%             A_{i+3/2} J_{i+3/2} - A_{i-1/2} J_{i-1/2} ≈ 2*(A_{i+1} J_{i+1} - A_i J_i)
%           (where A_i and J_i are evaluated at the mesh-point r_i = r_{i+1/2})
%           to simplify the boundary-current terms.
%
% Step 3 — Discretize the Fick's law (1st-moment) equation to get J_i in
%           terms of the cell-edge fluxes. Two options appear in the notes:
%             (a) Spherical (volume-weighted): J_i = -D_i * (A_i/V_i) * (phi_{i+1/2} - phi_{i-1/2})
%             (b) Cartesian (simple): J_i = -D_i / h_i * (phi_{i+1/2} - phi_{i-1/2})
%           where D_i = 1/(3*sigma_t,i) and h_i = r_{i+1/2} - r_{i-1/2}.
%           Think about which form is "nearly consistent" with the transport sweep.
%
% Step 4 — Substitute the expression for J_i from Step 3 into the combined
%           balance equation from Step 2.  The result should be a tridiagonal
%           system in the cell-edge unknowns phi_{i+1/2}.
%
% Step 5 — Boundary conditions:
%           (a) Inner boundary (r=0, reflective): J_{1/2} = 0 by symmetry.
%               Write out the first row of the tridiagonal system using this.
%           (b) Outer boundary (r=R, vacuum): Apply the Marshak (Mark) condition
%               relating the incoming partial current to the flux, i.e.
%               J_{N+1/2} = (1/4)*phi_{N+1/2}  (incoming partial current = 0).
%               Substitute into the last equation of the system.
%
% Step 6 — Assemble the final matrix equation A*phi = b, identify the
%           three diagonals, and verify it has the standard diffusion form.

\end{document}

